\documentclass[11pt,a4paper]{article}

\usepackage{fontspec}
\setmainfont{DejaVu Serif}[
  BoldFont={DejaVu Serif Bold},
  ItalicFont={DejaVu Serif},
  ItalicFeatures={FakeSlant=0.18},
  BoldItalicFont={DejaVu Serif Bold},
  BoldItalicFeatures={FakeSlant=0.18}
]
\setsansfont{DejaVu Sans}
\usepackage{polyglossia}
\setmainlanguage{english}
\usepackage[a4paper,margin=20.5mm]{geometry}
\usepackage{microtype}
\usepackage{setspace}
\setstretch{1.0}
\usepackage{amsmath,amssymb,amsthm,mathtools,bm}
\usepackage{booktabs,longtable,array,tabularx}
\usepackage{enumitem}
\usepackage{xcolor}
\usepackage[unicode,hidelinks]{hyperref}
\usepackage[nameinlink,noabbrev]{cleveref}
\emergencystretch=3em
\allowdisplaybreaks

\hypersetup{
  pdftitle={Bobsleigh Contact Geometry as a Typed Observation System},
  pdfauthor={Stas, Independent Research Program},
  pdfsubject={Mass geometry, unilateral contact, frictional dissipation, and observation loss},
  pdfkeywords={bobsleigh, contact mechanics, Lagrange-d'Alembert, friction, dissipation, observation geometry}
}

\definecolor{obsblue}{RGB}{37,83,138}
\definecolor{obsred}{RGB}{150,45,45}
\definecolor{lightblue}{RGB}{239,246,253}
\definecolor{lightred}{RGB}{253,242,242}

\theoremstyle{definition}
\newtheorem{definition}{Definition}[section]
\newtheorem{model}[definition]{Model}
\theoremstyle{plain}
\newtheorem{theorem}[definition]{Theorem}
\newtheorem{proposition}[definition]{Proposition}
\newtheorem{corollary}[definition]{Corollary}
\theoremstyle{remark}
\newtheorem{remark}[definition]{Remark}
\newtheorem{warning}[definition]{Boundary}

\newcommand{\R}{\mathbb{R}}
\newcommand{\dd}{\mathop{}\!d}
\newcommand{\norm}[1]{\left\lVert #1\right\rVert}
\newcommand{\pair}[2]{\left\langle #1,#2\right\rangle}
\newcommand{\SO}{\mathrm{SO}}
\newcommand{\SE}{\mathrm{SE}}
\newcommand{\Ocal}{\mathcal O}
\newcommand{\Kcal}{\mathcal K}
\newcommand{\Rcal}{\mathcal R}
\newcommand{\Acal}{\mathcal A}
\newcolumntype{L}[1]{>{\raggedright\arraybackslash}p{#1}}
\newcolumntype{Y}{>{\raggedright\arraybackslash}X}

\title{\bfseries Bobsleigh Contact Geometry\\
\large Mass distribution, unilateral contact, and typed observation loss}
\author{Stas\\
\small Independent Research Program; ORCID: 0009-0000-2294-705X}
\date{Strict GO Core reference revision v1.1 --- 28 July 2026}

\begin{document}
\maketitle

\begin{center}
\fcolorbox{obsblue}{lightblue}{%
\parbox{0.92\linewidth}{\small
\textbf{Status.} This document supersedes the bilingual working draft
v1.0. It is migrated to GO Core v0.2 and the Frames--Forces--Dissipation
Interface v0.1. Contact and friction are diagnostic smooth-regime
models. No sled design, ice constitutive law, safety calculation, or
engineering validation is claimed.}}
\end{center}

\begin{abstract}
A center-of-mass trace does not determine the contact state of a sled
moving on a shaped surface. We formalize this observation using a
regular track parametrization, a finite mass realization, a
configuration-space mass form, unilateral gaps, normal reactions, and
tangential contact forces. The mass form
\[
G_{ij}(q)=\sum_a m_a\,\partial_i x_a(q)\cdot\partial_jx_a(q)
\]
is positive semidefinite and is a metric precisely where the
realization map is an immersion. Contact uses the fixed convention
\[
g_j\ge0,\qquad N_j\ge0,\qquad N_jg_j=0,\qquad
\norm{T_j}\le\mu_jN_j .
\]
With \(T_j\) defined as the tangential force on the sled and
\(v_{T,j}\) as its velocity relative to the ice, the nonnegative
frictional loss is
\[
P_{\rm fric}=-\sum_jT_j\cdot v_{T,j}\ge0.
\]
Consequently the autonomous stationary-track energy balance is
\(\dot E=P_{\rm act}-P_{\rm fric}\) in the smooth active-contact
regime. This corrects the sign in v1.0. Cone utilization is defined
piecewise at zero normal load, and load entropy has an explicit
logarithm base and normalization. Center-of-mass, pose, inertial,
contact, load, and thermal channels remain distinct observations of a
larger hidden state.
\end{abstract}

\tableofcontents

\section{Scope and protocol}

Let \(\lambda\) specify a track chart, world and sled frames, clock,
contact convention, sampling rates, sensor calibration, estimator, and
uncertainty model. The observation chain is
\[
X_t\xrightarrow{\Rcal}\widetilde X_t
\xrightarrow{\Ocal_\lambda}Z_t
\xrightarrow{\Kcal_\lambda}\operatorname{Law}(Y_t\mid Z_t)
\xrightarrow{\widehat\theta_\lambda}\widehat\theta_t.
\]
A kinematic centerline supplied as input is not a dynamically
predicted trajectory. A proxy load rule is not a solved contact
reaction. These statuses are recorded separately.

\begin{warning}[Smooth-regime boundary]
The equations below apply between impacts while the active contact set
is fixed and the state is differentiable. Impact impulses,
restitution, Painlev\'e-type inconsistency, stick--slip switching, and
ice deformation require a nonsmooth differential inclusion or
complementarity time-stepper and are not proved here.
\end{warning}

\section{Track as a shaped surface}

Let \((e_s,e_u,n_{\rm tr})\) be an orthonormal track frame and let
\(\gamma(s)\) be a reference curve. A local trough chart is
\begin{equation}\label{eq:surface}
X_{\rm tr}(s,u)
=\gamma(s)+R_{\rm tr}(s)
\begin{pmatrix}0\\u\\h(s,u)\end{pmatrix},
\qquad R_{\rm tr}=(e_s,e_u,n_{\rm tr})\in\SO(3).
\end{equation}
Its admissible domain \(D\subset\R^2\) must satisfy
\begin{equation}\label{eq:surface-regular}
\partial_sX_{\rm tr}\times\partial_uX_{\rm tr}\ne0
\quad\text{on }D.
\end{equation}
The surface normal is therefore
\[
n_\Sigma
=\frac{\partial_sX_{\rm tr}\times\partial_uX_{\rm tr}}
{\norm{\partial_sX_{\rm tr}\times\partial_uX_{\rm tr}}}.
\]
Centerline curvature, banking, transverse trough curvature, and chart
regularity are distinct data. A centerline alone does not determine
contact geometry.

\section{Finite mass realization}

Let a rigid sled be represented by masses \(m_a>0\) at body points
\(r_a^B\). Its total mass, center of mass, and inertia about the center
of mass are
\begin{align}
M&=\sum_{a=1}^nm_a,&
r_{\rm cm}^B&=\frac1M\sum_{a=1}^nm_ar_a^B,\\
I_B&=\sum_{a=1}^nm_a
\left(
\norm{\rho_a}^2I_3-\rho_a\rho_a^{\mathsf T}
\right),&
\rho_a&=r_a^B-r_{\rm cm}^B .
\end{align}
For world origin \(x^I(t)\) and attitude
\(A_{I\leftarrow B}(t)\in\SO(3)\),
\[
x_a^I=x^I+A_{I\leftarrow B}\rho_a.
\]

Let \(q\in Q\) be local generalized coordinates after a declared
reduction, and let
\(\Phi(q)=(x_1(q),\ldots,x_n(q))\). Define
\begin{equation}\label{eq:mass-form}
G_{ij}(q)=\sum_{a=1}^nm_a
\partial_i x_a(q)\cdot\partial_jx_a(q).
\end{equation}

\begin{proposition}[When the mass form is a metric]
\(G(q)\) is positive semidefinite. It is positive definite if and only
if \(\dd\Phi_q\) is injective.
\end{proposition}

\begin{proof}
For \(\xi\in T_qQ\),
\[
\xi^{\mathsf T}G(q)\xi
=\sum_am_a\norm{\dd\Phi_q(\xi)_a}^2.
\]
\end{proof}

Thus redundant coordinates produce a degenerate mass form. Calling it
a Riemannian metric requires the immersion hypothesis.

\section{Unilateral contact and force pullback}

For candidate contact \(j\), let \(g_j(t,q)\) be the signed gap, with
\(g_j>0\) meaning separation. In the smooth force regime:
\begin{equation}\label{eq:comp}
g_j\ge0,\qquad N_j\ge0,\qquad N_jg_j=0.
\end{equation}
Let \(n_j^I\) be the ice-to-sled unit normal and
\(P_{T,j}=I-n_j^I(n_j^I)^{\mathsf T}\). The contact force on the sled
is
\begin{equation}\label{eq:contact-force}
F_j^I=N_jn_j^I+T_j^I,\qquad
T_j^I\cdot n_j^I=0.
\end{equation}
If \(J_j^I=\partial_qc_j^I\) is the contact-point Jacobian, its
generalized force is
\[
Q_j=J_j^{\mathsf T}F_j^I\in T_q^*Q.
\]
This pullback is required before adding \(Q_j\) to other generalized
forces.

Let
\[
v_{T,j}^I
=P_{T,j}\bigl(v_{{\rm sled},j}^I-v_{{\rm ice},j}^I\bigr).
\]
A Coulomb-type admissibility condition is
\begin{equation}\label{eq:friction-cone}
\norm{T_j^I}\le\mu_jN_j,\qquad
T_j^I\cdot v_{T,j}^I\le0.
\end{equation}
For sliding, a maximum-dissipation selection is
\[
T_j^I=-\mu_jN_j
\frac{v_{T,j}^I}{\norm{v_{T,j}^I}}.
\]
For sticking, \(v_{T,j}=0\) and \(T_j\) is set-valued inside the cone.

\section{Dynamics and corrected energy balance}

In a smooth reduced chart, let
\[
L(t,q,\dot q)
=\frac12\dot q^{\mathsf T}G(t,q)\dot q-V(t,q).
\]
The Lagrange--d'Alembert equation is
\begin{equation}\label{eq:lda-bob}
\frac{\dd}{\dd t}\partial_{\dot q}L-\partial_qL
=Q_{\rm act}+Q_{\rm fric}+Q_{\rm con}.
\end{equation}
With
\[
E_L=\partial_{\dot q}L\cdot\dot q-L,
\]
the frame--force interface gives
\begin{equation}\label{eq:energy-general}
\dot E_L
=P_{\rm act}-P_{\rm fric}
+P_{\rm con}-\partial_tL,
\end{equation}
where the nonnegative frictional dissipation is
\begin{equation}\label{eq:pfric}
\boxed{
P_{\rm fric}
=-\sum_jT_j^I\cdot v_{T,j}^I\ge0 .
}
\end{equation}

\begin{corollary}[Stationary ideal support]
If the track is stationary, normal constraints are ideal, \(L\) is
autonomous, and no other nonconservative forces act, then between
impacts
\[
\dot E_L=-P_{\rm fric}\le0.
\]
\end{corollary}

\begin{remark}[Correction to v1.0]
Version v1.0 defined \(T_j\) as the tangential force on the sled but
wrote
\(\dot E=-P_{\rm fric}\) with
\(P_{\rm fric}=\sum_jT_j\cdot v_{j,\rm rel}\).
Under dissipative friction that sum is nonpositive, so the displayed
energy law increased energy. Equation \cref{eq:pfric} defines the loss
with the required leading minus sign.
\end{remark}

For a moving support, normal reactions can exchange power and
\(P_{\rm con}\) need not vanish. The same support work must not also
be inserted through a reduced time-dependent \(L\).

\section{Guarded diagnostics}

\subsection{Cone utilization}

Define
\[
u_j=
\begin{cases}
\norm{T_j}/(\mu_jN_j),&\mu_jN_j>0,\\
0,&\mu_jN_j=0\ \text{and}\ T_j=0 .
\end{cases}
\]
The cone law excludes \(\mu_jN_j=0,T_j\ne0\). For a finite candidate
set \(J\),
\[
S_{\rm slip}=1-\max_{j\in J}u_j,\qquad
\max\varnothing:=0.
\]
Then \(0\le S_{\rm slip}\le1\) on admissible states. This is a cone
margin, not a probability of slipping.

\subsection{Load distribution and entropy}

If \(N_{\rm tot}=\sum_{j=1}^KN_j>0\), set
\[
p_j=\frac{N_j}{N_{\rm tot}},\qquad
H_b(N)=-\sum_{j:p_j>0}p_j\log_b p_j,
\quad b>1.
\]
For a declared \(K\ge2\), the normalized diagnostic is
\[
\overline H_b=\frac{H_b(N)}{\log_bK}\in[0,1].
\]
It is independent of the logarithm base. If \(N_{\rm tot}=0\), the
load distribution and entropy are undefined rather than silently set
to zero.

\subsection{Two-runner imbalance}

For \(N_L+N_R>0\),
\[
\Delta_N
=\frac{\lvert N_R-N_L\rvert}{N_R+N_L}\in[0,1].
\]
The diagnostic is undefined in free flight.

\subsection{Geometric demand and banking mismatch}

For a centerline with curvature \(\kappa\) and speed \(v\),
\[
a_\kappa=v^2\kappa
\]
has acceleration dimension. If \(f_{\rm req}=m(a-g)\ne0\), define the
unsigned mismatch
\[
\Delta_b
=\operatorname{atan2}
\left(
\norm{n_{\rm tr}\times f_{\rm req}},
n_{\rm tr}\cdot f_{\rm req}
\right)\in[0,\pi].
\]
This is a directional diagnostic; it does not solve the contact
forces.

\section{Observation channels}

The hidden record may contain
\[
X_t=(\Sigma,\{m_a,r_a^B\},q,\dot q,\Acal_t,
\{c_j,N_j,T_j\}_{j\in\Acal_t},\vartheta),
\]
where \(\Acal_t\) is the active contact set and \(\vartheta\) collects
constitutive parameters. Useful deterministic channels include
\begin{align*}
\Ocal_{\rm cm}(X_t)&=x_{\rm cm}^I(t),&
\Ocal_{\rm pose}(X_t)&=A_{I\leftarrow B}(t),\\
\Ocal_{\rm imu}(X_t)&=(f^B,\Omega^B),&
\Ocal_{\rm load}(X_t)&=(N_j)_{j\in\Acal_t},\\
\Ocal_{\rm contact}(X_t)&=(c_j)_{j\in\Acal_t},&
\Ocal_{\rm heat}(X_t)&=P_{\rm fric}(t).
\end{align*}
No equality between these channels is implied.

\begin{proposition}[Conditional center-of-mass non-identifiability]
Assume the admissible hidden-state class contains two records
\(X_t^{(1)}\ne X_t^{(2)}\) with the same center-of-mass curve but
different mass realization or contact record. Then
\(\Ocal_{\rm cm}\) is not injective on that class.
\end{proposition}

\begin{proof}
By hypothesis,
\(\Ocal_{\rm cm}(X_t^{(1)})=\Ocal_{\rm cm}(X_t^{(2)})\) while the
arguments differ.
\end{proof}

This proposition is deliberately conditional. Showing that two
specific sleds admit the same trajectory under a fixed constitutive
model is a separate dynamical existence problem.

\section{Reproducibility record}

A numerical benchmark is admissible only if it reports:
\begin{enumerate}[label=(\roman*)]
  \item the surface chart and its regularity domain;
  \item mass points, center of mass, inertia, and coordinate chart;
  \item active-set and contact-law solver;
  \item ice-motion convention and friction parameters;
  \item integration method, tolerances, event handling, and convergence
  study;
  \item sensor map, sampling, noise, and estimator;
  \item which outputs are solved reactions and which are proxies.
\end{enumerate}
The proxy values in v1.0 are not retained as validated mechanics
because no complementarity/contact solver or convergence study was
attached.

\section{Claim boundary}

The strict result is a corrected and typed interface to standard
constrained mechanics. It does not claim:
\begin{itemize}
  \item a predictive ice--runner friction law;
  \item uniqueness of the contact solution;
  \item engineering accuracy or safety certification;
  \item physiological rider limits;
  \item a new theorem beyond the stated interface propositions.
\end{itemize}

\begingroup
\small
\hbadness=10000
\begin{thebibliography}{9}

\bibitem{brogliato}
B. Brogliato, \emph{Nonsmooth Mechanics: Models, Dynamics and
Control}, 3rd ed., Springer, 2016.
\href{https://doi.org/10.1007/978-3-319-28664-8}
{doi:10.1007/978-3-319-28664-8}.

\bibitem{bloch}
A.~M. Bloch, \emph{Nonholonomic Mechanics and Control}, 2nd ed.,
Springer, 2015.
\href{https://doi.org/10.1007/978-1-4939-3017-3}
{doi:10.1007/978-1-4939-3017-3}.

\bibitem{acary}
V. Acary and B. Brogliato,
\emph{Numerical Methods for Nonsmooth Dynamical Systems},
Springer, 2008.
\href{https://doi.org/10.1007/978-3-540-75392-6}
{doi:10.1007/978-3-540-75392-6}.

\end{thebibliography}
\endgroup

\end{document}
